\chapter{Formaliser les mathématiques}
\label{chp.axiomes}

\minitoc

\lettrine{P}{our} commencer notre étude, nous allons étudier l'une des
propriétés les plus connues de la théorie des ensembles~: la capacité
d'expression suffisante pour formaliser toutes les mathématiques usuelles. Ce
chapitre se concentrera donc sur l'étude des axiomes de ZFC et de leur
utilisation pour construire les objets mathémtiques que nous connaissons
habituellement.

Nous devons ici faire quelques éclaircissements d'ordre philosophique. Tout
d'abord, sur la méta-théorie~: nous considérons que l'univers ambiant dans
lequel nous pratiquons les mathématiques est lui-même un modèle de la théorie
ZFC. Plus précisément, en notant $\unZF$ l'objet mathématique qui sera la
collection de tous les objets mathématiques, ZFC peut être vu comme une théorie
donnant une approximation du comportement de $\unZF$.

En particulier, $\unZF$ contient des collections, que nous appelons
d'habitude ensembles, et si un objet $X$ apparait dans une collection $\varCC$,
on écrit $X\in\varCC$. Cependant, les objets que nous manipulerons ne seront pas
ces collections, et la relation $\in$ utilisée dans $\ZFC$ ne sera pas \og être
un objet mathématique apparaissant dans cette collection intuitive\fg{}~: nous
utiliserons des outils purement formels pour étudier $\ZFC$, et devons donc les
distinguer des objets intuitifs associés. Aussi nous appellerons \og ensemble\fg
un objet de la théorie $\ZFC$, et \og collection\fg un ensemble au sens
intuitif, dans $\unZF$. De même, les objets formels seront associés à la notion
d'appartenance, et les collections à la notion d'occurrence.

Remarquons cependant qu'un ensemble donne de fait naissance à une
collection~: l'ensemble $X$ définit la collection
$\compr{ x \in\unZF}{x \in X}$.

Enfin, nous traiterons régulièrement des classes. S'il est possible de ne se
restreindre qu'aux ensembles pour l'étude de $\ZFC$, nous verrons qu'il est bien
plus pratique d'énoncer certaines constructions en terme de classes. Une classe
est une collection qui ne peut pas être représentée par un ensemble (par exemple
la classe des ensembles, \latinexpr{cf}. le paradoxe de Russell), mais que l'on
peut décrire par un prédicat. Ainsi, une classe peut se représenter par le
prédicat lui correspondant. Si une classe $\varCC$ est décrite par un prédicat
$\varFF(x)$, alors $x\in\varCC$ doit se lire comme une reformulation plus
lisible de $\varFF(x)$. Un léger défaut vient avec cette approche~: certains
théorèmes, se référant aux classes, doivent se lire comme des schémas de
théorèmes, c'est-à-dire des théorèmes paramétrés par le prédicat $\varFF$
définissant une classe.

\section{Axiomes de ZFC}

\subsection[Premiers axiomes]{Extensionalité, paire, union, ensemble de parties}

Commençons par étudier la théorie $\ZFC$ en en listant les axiomes. Un premier
point important est de définir le langage sur lequel nous travaillerons~: nous
utiliserons $\SignP{\ZF} = \{ \in^2\}$. En effet, tous les énoncés seront écrits
seulement à l'aide du symbole de relation binaire $\in$ et de la relation $=$,
symbole logique déjà inclus dans tout langage.

\begin{remark}
  Si, techniquement, la théorie ensembliste est une théorie relationnelle, nous
  utilisons abondamment des extensions par définition
  (\latinexpr{c.f.} \cref{def.ext.def.rel} et
  \cref{def.ext.def.fun}) pour introduire de nouveaux symboles~:
  $\{x,y\}, \bigcup x, x \subseteq y,$ \latinexpr{etc.}
  Il est néanmoins important de ne pas
  les inclure dans notre signature en première instance, car le sens des
  nouveaux symboles est uniquement déterminé par le sens de $\in$ et des axiomes
  de $\ZF$. Ainsi, dans le \cref{chp.modZFC}, nous considérons des
  sous-modèles (en un sens large) de l'univers $\unZF$ mais uniquement vis
  à vis de la relation $\in$, les autres opérations pouvant différer d'une
  sous-modèle à un autre.
\end{remark}

\begin{axiom}[Extensionalité]\label{ax.ZF.ext}
  L'axiome d'extensionalité exprime que deux ensembles sont égaux exactement
  lorsqu'ils possèdent les mêmes éléments~:
  \[\forall x \forall y, \qquad x = y \iff (\forall z, z \in x\iff z \in y)\]
\end{axiom}

Cet axiome définit ce que signifie l'égalité dans le monde des ensembles. Cela
permet directement de voir que l'ordre ou le nombre d'occurrences n'importent
pas dans un ensemble, contrairement par exemple au cas des listes. C'est un
axiome au statut particulier car il ne permet pas de construire de nouvel
ensemble. Les autres axiomes, eux, donneront principalement de nouvelles
méthodes pour, à partir d'ensembles déjà construits, définir de nouveaux
ensembles.

\begin{axiom}[Paire]\label{ax.ZF.pair}
  L'axiome de la paire exprime que si deux ensembles $x$ et $y$ ont été
  construits, alors l'ensemble $\{x,y\}$ peut être construit à partir d'eux~:
  \[\forall x \forall y \exists z, \qquad\forall a, a \in z \iff
  a = x \lor a = y\]
\end{axiom}

\begin{exercise}
  Montrer que pour tous ensembles $x,y$, il existe un unique ensemble $z$ tel
  que décrit dans l'\cref{ax.ZF.pair}. Cela justifie donc la notation $\{x,y\}$
  pour cet ensemble.
\end{exercise}

En utilisant l'axiome de la paire, on peut ainsi construire des collections
finies telles que $\{x,\{y,z\}\}$, mais on ne peut pas par exemple définir
$\{x,y,z\}$~: il nous faut un axiome permettant d'accéder à l'intérieur d'un
ensemble pour en construire un nouveau.

\begin{axiom}[Union]\label{ax.ZF.union}
  Pour tout ensemble $x$, on peut construire l'ensemble $\bigcup x$ contenant
  les éléments des éléments de $x$~:
  \[\forall x \exists y,\qquad \forall z, z\in y \iff
  (\exists a, z \in a \land a \in x)\]
\end{axiom}

\begin{notation}
  De la même façon qu'avec l'axiome de la paire, l'ensemble $y$ défini par
  l'axiome est unique, on l'écrira donc $\bigcup x$.
\end{notation}

\begin{exercise}
  Soient $x_1,\ldots,x_n$ des ensembles, montrer par récurrence qu'il existe
  l'ensemble $\{x_1,\ldots,x_n\}$.
\end{exercise}

\begin{notation}
  Etant donnée une famille $x_1,\ldots,x_n$ d'ensembles, on définit
  $\displaystyle \bigcup_{i = 1}^n x_i$ comme $\bigcup\{x_1,\ldots,x_n\}$. En
  particulier, $x \cup y$ est $\bigcup\{x,y\}$.
\end{notation}

Enfin, l'axiome de l'ensemble des parties permet de considérer comme un ensemble
la collection des parties d'un ensemble. En un sens, il permet de faire croître
considérablement la taille de ce que l'on peut construire.

\begin{notation}
  On définit le prédicat binaire $\subseteq$, appelé l'inclusion, par
  \[x \subseteq y \defeq \forall z, z \in x \implies z \in y\]
\end{notation}

\begin{axiom}[Ensemble des parties]\label{ax.ZF.pow}
  Pour tout ensemble $x$, il existe l'ensemble $\powerset (x)$ dont les éléments
  sont exactement les ensembles inclus dans $x$~:
  \[\forall x\exists y, \qquad \forall z, z \in y \iff z \subseteq x\]
\end{axiom}

\begin{notation}
  Comme le $y$ défini plus haut est unique, on peut là encore lui donner une
  notation, qui est bien sûr $\powerset(x)$.
\end{notation}

\subsection{Les schémas d'axiomes}

Les axiomes précédemment donnés constituent les briques de base pour construire
des ensembles, mais sont en général trop grossiers. L'intérêt de la théorie des
ensembles est de pouvoir construire des ensembles correspondant à des
collections, ce qui manque pour l'instant à notre système. C'est pour cela que
nous ajoutons l'axiome de compréhension~: étant donnée une formule $\varFF$ et
un ensemble $X$, on peut construire l'ensemble $\compr{x \in X}{\varFF(x)}$ des
éléments de $x$ vérifiant $\varFF$.

La restriction de la compréhension à des parties d'un ensemble s'explique par le
paradoxe de Russell~: si l'on pouvait construire un ensemble pour chaque
formule, on pourrait construire $\compr{x}{x\notin x}$, qui appartient à
lui-même si et seulement s'il n'appartient pas à lui-même.

Un autre problème doit être contourné, et il est la raison pour laquelle on
parle de schéma d'axiomes et non d'axiome. Si l'on voulait définir le schéma
avec ce qui a été dit, celui-ci commencerait moralement par $\forall \varFF$~:
cela n'est pas une quantification du premier ordre, et il n'est donc pas
possible de donner un axiome pour toutes les formules. \'A la place, on
introduit un schéma d'axiomes, c'est-à-dire une infinité d'axiomes ayant tous
la même forme et dépendant d'un paramètre que l'on quantifie sur les formules.

\begin{axiom}[Schéma de compréhension]\label{ax.ZF.compre}
  Soit $\varFF(x_0,\ldots,x_n)$ une formule sur $\SignP{\ZF}$
  à $n$ variables libres. Alors pour tous ensembles $X$ et $a_1,\ldots,a_n$, il
  existe l'ensemble $\compr{x\in X}{\varFF(x,a_1,\ldots,a_n)}$, ce qui s'écrit
  formellement
  \[\forall X\forall a_1\;\cdots\;a_n\exists y,\qquad
    \forall x, x \in y \iff (x\in X\land \varFF(x,a_1,\ldots,a_n))\]
\end{axiom}

\begin{notation}
  On définit donc la notation $\compr{x\in X}{\varFF(x,a_1,\ldots,a_n)}$ pour
  l'ensemble précédent.
\end{notation}

\begin{exercise}
  Soit un ensemble $x$ non vide, montrer qu'il existe l'ensemble $\bigcap x$ des
  éléments qui sont dans tous les éléments de $x$.
\end{exercise}

\begin{notation}
  On définit des notations pour l'intersection analogues à celles pour l'union~:
  $\displaystyle\bigcap_{i = 1}^n x_i$ et $x \cap y$.
\end{notation}

\begin{exercise}\label{exo.ZF.prod}
  Soient $x$ et $y$ deux ensembles. On définit le couple $(x,y)$ par
  \[(x,y) \defeq \{\{x,y\},\{x\}\}\]
  Montrer que pour tous $x,y,x',y'$, $(x,y) = (x',y')$ si et seulement si
  $x=x'$ et $y=y'$.

  Construire un prédicat $\varFF(x,y,z)$ tel que $\varFF(x,y,z)$ est vrai
  si et seulement si $z = (x,y)$. En déduire en considérant une partie bien
  choisie de $\powerset (\powerset (\powerset(x\cup y)))$ que
  \[x\times y \defeq \compr{(a,b)}{a \in x, b \in y}\]
  est un ensemble bien défini.
\end{exercise}

Le second schéma, le schéma d'axiomes de remplacement, peut être vu comme une
version plus forte de la compréhension. Plutôt que de s'intéresser à filtrer des
éléments dans un ensemble plus gros, le but de ce schéma d'axiomes est de
construire un ensemble par une fonction. Comme la notion de fonction n'est pas
encore définie, nous utilisons à la place la notion de relation fonctionnelle.

\begin{definition}[Relation fonctionnelle]
  Une relation binaire est ici une formule à deux variables libres. On dit
  qu'une relation $\varRR$ est fonctionnelle si pour chaque $x$, il existe
  au plus un $y$ tel que $x \varRR y$. On écrira pour raccourcir
  \[\Funct(\varRR) \defeq \forall x\forall y\forall z,
  x \varRR y\land x \varRR z \implies y = z\]

  Pour une relation fonctionnelle $\varRR$, on définit la collection du
  domaine de $\varRR$~:
  \[\dom(\varRR)(x)\defeq \exists y, x \varRR y\]
  et la collection de l'image de $\varRR$~:
  \[\im(\varRR)(y)\defeq \exists x, x \varRR y\]
\end{definition}

\begin{axiom}[Schéma de remplacement]\label{ax.ZF.repl}
  Pour toute formule $\varFF(x_0,\ldots,x_{n+1})$, pour tous ensembles
  $X,a_1,\ldots,a_n$, si $\varFF(a_1,\ldots,a_n)$ est une relation fonctionnelle
  en ses deux dernières variables libres alors l'ensemble image de
  $\varFF(a_1,\ldots,a_n)$ par $X$ est aussi un ensemble~:
  \begin{multline*}
    \forall a_1\;\cdots\;a_n, \Funct(\varFF(a_1,\ldots,a_n))\implies \forall X,
    \exists y,\\
    \forall x, x \in y \iff (\exists z, z\in X \land \varFF(a_1,\ldots,a_n,z,x))
  \end{multline*}
\end{axiom}

\begin{exercise}
  En remarquant qu'une relation fonctionnelle peut représenter une fonction
  partielle, montrer que le schéma d'axiomes de compréhension peut se déduire du
  schéma d'axiomes de remplacement.
\end{exercise}

\begin{remark}
  Le schéma de remplacement peut être décomposé en deux schémas plus faibles~:
  \begin{itemize}
  \item le schéma de compréhension
  \item le schéma de remplacement total, énonçant la même chose que le schéma
    de remplacement mais demandant à la relation fonctionnelle d'être totale,
    c'est-à-dire de vérifier le prédicat
    \[\FunTot(\varRR) \defeq \forall x \exists! y, x \varRR y\]
    plutôt que $\Funct$.
  \end{itemize}

  \noindent
  En pratique, on utilise plutôt ces deux schémas indépendamment. Dans le cas
  où on a défini une fonction $f$ et pour $X$ un ensemble, on notera alors
  \[\compr{f(x)}{x \in X}\]
  l'ensemble construit par cet axiome qui est défini de façon unique. Dire
  qu'un objet $y$ appartient à cet ensemble revient donc à dire qu'il existe
  $x \in X$ tel que $y = f(x)$.
\end{remark}

\begin{exercise}\label{exo.ZF.prod2}
  En réutilisant le prédicat $\varFF$ de l'\cref{exo.ZF.prod}, montrer grâce au
  schéma d'axiomes de remplacement que pour tous ensembles $x,y$, l'ensemble
  $x\times y$ est bien défini même sans l'axiome de l'ensemble des parties.
\end{exercise}

\noindent
Donnons aussi l'axiome de l'univers non vide.

\begin{axiom}[Univers non vide]\label{ax.ZF.nonEmpty}
  Il existe un ensemble~:
  \[\exists x, x = x\]
\end{axiom}

\begin{exercise}
  Montrer que l'axiome précédent est équivalent à l'existence de l'ensemble vide
  $\vide$ défini par
  \[\forall z, z\notin \vide\]
  (où $x\notin y$ signifie $\lnot(x\in y)$) 
\end{exercise}

\begin{remark}
  Le fait d'ajouter ou non cet axiome dépend de si l'on considère qu'un modèle
  est nécessairement non vide (de façon équivalente, que l'énoncé
  $(\forall x, \varFF(x)) \implies (\exists x,\varFF(x))$ est vérifié). Avec
  la syntaxe de preuves donnée dans le \cref{chp.logpred}, cet axiome n'est pas
  nécessaire. Nous verrons de toute façon que l'axiome de l'infini nous donnera
  automatiquement un ensemble dans notre univers, mais nous présentons cet
  axiome pour pouvoir considérer la théorie sans axiome de l'infini mais où
  l'univers est quand même supposé non vide.
\end{remark}

\subsection{L'axiome de l'infini et les entiers}

Pour introduire l'axiome suivant, il nous faut d'abord motiver l'utilisation de
ses éléments constitutifs. L'ensemble mathématique le plus élémentaire que l'on
est amené à étudier est certainement $\bN$, l'ensemble des entiers
naturels. Une formalisation habituelle de cet ensemble demande en général les
trois constituants suivant~:
\begin{itemize}
\item l'élément $0$
\item la fonction unaire $\Ss$ (successeur), correspondant à $n \mapsto n + 1$
\item le principe de récurrence, que l'on peut encoder dans les ensembles par le
  fait que si $\varPt\subseteq\bN$, $0 \in \varPt$ et
  $\forall n, n\in \varPt \implies \Ss n \in \varPt$ alors $\varPt = \bN$.
\end{itemize}

Chercher à définir $\bN$ dans $\ZFC$ nous demande donc de définir ces
éléments. Un candidat naturel à $0$ est $\vide$~: les deux sont les objets
nuls par excellence, et $\vide$ est le premier ensemble que l'on peut
construire. La question de savoir ce qu'est $\Ss x$ pour un ensemble $x$ est
alors naturelle~: pour cela, nous utilisons le codage de Von Neumann consistant
à coder l'entier naturel $n$ par $\intervInt{n}$. Cet encodage nous offre
une définition simple à la fonction $\Ss$~:
\[\Ss x\defeq x\cup \{x\}\]

Remarquons que l'on peut déjà construire tous les entiers que l'on souhaite~:
on peut construire $0$ et itérer la fonction $\Ss$. Malheureusement, rien ne
nous dit que la collection $\compr{\Ss^n\;0}{n \in \bN}$ est bien elle-même
un ensemble (où le $\bN$ apparaissant dans la définition est l'ensemble
des entiers naturels de notre méta-théorie, nos entiers au sens intuitif). Pour
pallier ce problème, et pour éviter d'utiliser notre méta-théorie explicitement,
on va à la place définir $\bN$ comme le plus petit ensemble contenant $0$
et stable par la fonction $\Ss$. Il nous reste à savoir qu'un ensemble contenant
$0$ et stable par $\Ss$ existe bien.

\begin{axiom}[Infini]\label{ax.ZF.infini}
  Il existe un ensemble contenant $\vide$ et stable par $\Ss$~:
  \[\exists x, \qquad \vide \in x \land
  (\forall y, y \in x \implies \Ss y\in x)\]
\end{axiom}

\begin{remark}
  On peut donc se passer de l'axiome de l'ensemble vide en prenant l'axiome de
  l'infini (et en adaptant sa définition pour ne pas appeler explicitement
  l'ensemble vide). Donnons ainsi une version de l'axiome de l'infini ne
  mentionnant que la relation $\in$ et l'égalité~:
  \begin{multline*}
    \exists x,
    (\forall y, (\forall z, \lnot z \in y) \implies y \in x) \mathrel{\land}\\
    (\forall y, y \in x \implies
    (\forall z, (\forall a, a \in z \iff a = y \lor a \in y) \implies
    z \in x))
  \end{multline*}
  On reconnait les deux parties du $\land$~: à gauche, on dit que tout ensemble
  vide est dans $x$, et à droite que tout élément $y$ de $x$ contient tous les
  éléments de la forme $\Ss y$ (c'est-à-dire, tous les $z$ dont les éléments
  sont ceux de $y$ et $y$). Il est clair que cette définition est à la fois
  longue et difficile à lire, d'où l'intérêt d'introduire de nouveaux symboles.
\end{remark}

\begin{definition}[Entiers naturels]
  Soit $X$ l'ensemble défini par l'axiome de l'infini. On définit alors
  $\bN$ comme
  \[\bN \defeq \compr{x\in X}{\forall \varPt, \varPt\subseteq X\land
    \vide \in \varPt \land (\forall a, a \in \varPt \implies \Ss a\in \varPt)
    \implies x\in \varPt}\]
\end{definition}

\noindent
On vérifie maintenant le principe de récurrence.

\begin{theorem}[Récurrence]
  Soit $\varPt$ une partie de $\bN$ telle que $\vide \in \varPt$ et telle que
  $\forall n, n \in \varPt \implies \Ss n \in \varPt$, alors $\varPt = \bN$.
\end{theorem}

\begin{proof}
  Comme $\varPt$ est une partie de $\bN$, il nous suffit de montrer que
  $\bN \subseteq \varPt$. Par transitivité de l'inclusion,
  $\varPt\subseteq X$ pour $X$ l'ensemble à partir duquel $\bN$ a été défini.
  On sait de plus que
  \[x \in \bN \implies \forall Y, Y\subseteq X \land
  \vide \in Y \land (\forall a, a \in Y \implies \Ss a \in Y)
  \implies x\in Y\]
  d'où, en spécialisant $Y$ en $\varPt$, et sachant que $\vide\in \varPt$ et
  $\forall a, a \in \varPt \implies \Ss a \in \varPt$, il vient que
  \[x\in \bN \implies x \in \varPt\] ce qui est exactement
  $\bN\subseteq \varPt$, d'où le résultat.
\end{proof}

\begin{notation}
  A partir de maintenant, pour fluidifier l'écriture, on adoptera un style plus
  laxiste sur l'écriture des propositions. Par exemple on se permettra d'écrire
  $\forall x \in X, \varFP$ pour $\forall x, x\in X \implies \varFP$ et tous les
  légers abus de notations du même genre.
\end{notation}

\begin{remark}
  Pour que la définition de $\bN$ soit robuste, il convient de vérifier
  que le $X$ utilisé ne change pas le résultat final~: c'est le cas. En effet,
  si deux ensembles $X$ et $Y$ contiennent $\vide$ et sont stables par
  successeur, alors notons $\bN_X$ et $\bN_Y$ les deux versions de
  $\bN$ obtenues respectivement avec $X$ et avec $Y$. Alors $X\cap Y$
  contient $\vide$ et est stable par successeur, donc $X\cap Y$ contient
  $\bN_X$ et $\bN_Y$, mais alors on en déduit que $\bN_Y$ est
  une partie de $X$ contenant $\vide$ et stable par successeur, donc
  $\bN_X \subseteq \bN_Y$. De même $\bN_Y\subseteq\bN_X$ donc, finalement,
  $\bN_X = \bN_Y$.
\end{remark}

\subsection{Axiome du choix et fonctions}

Pour introduire l'axiome du choix, le mieux est de parler de ce qu'on voudrait
faire mais ne peut pas faire sans lui. Prenons $X_1,\ldots,X_n$ des ensembles
finis non vides. On peut prouver qu'il existe
$(x_1,\ldots,x_n)\in X_1\times\cdots\times X_n$ (en codant le produit itéré et
les tuples de taille arbitraire). Pour cela, il suffit de raisonner par
récurrence sur $n$~: pour passer de l'étape $n$ à l'étape $n+1$ il suffit de
remarquer que $X_{n+1}$ est non vide, et prendre alors un élément $x\in X_{n+1}$
pour l'ajouter au tuple déjà construit.

Le problème ici est que, si l'on peut obtenir ce résultat sur un nombre fini
d'ensembles non vides, ce processus ne fonctionne plus lorsque l'on passe au cas
infini. On ne peut pas faire une infinité de choix, puisque l'on n'utilise que
des méthodes finies. C'est pourquoi l'axiome du choix existe~: il permet de
choisir dans des familles infinies d'ensembles non vides. Donnons d'abord la
définition de cet axiome.

\begin{axiom}[Choix]\label{ax.ZF.AC}
  Soit $X$ un ensemble dont les éléments sont non vides et deux à deux
  disjoints. Alors il existe un ensemble $\varCh$ tel que pour tout $x\in X$,
  $\varCh\cap x$ est un singleton. \'Ecrit formellement~:
  \[\forall X, \vide\notin X \land
  (\forall x,y \in X, x \cap y \neq
  \vide \implies x = y)
  \implies \exists \varCh, \forall x \in X, \exists y, \varCh \cap x = \{y\}\]
\end{axiom}

\begin{remark}
  En l'état, l'ensemble $\varCh$ obtenu par l'axiome du choix peut tout à fait
  contenir des éléments qui ne sont pas éléments d'éléments de $X$~: on impose
  simplement que les éléments d'éléments de $X$ qui sont dans $\varCh$ donnent
  la propriété souhaitée. On peut cependant prendre
  $\varCh'\defeq\varCh\cap \bigcup X$ pour avoir en plus le résultat
  \[\forall y, y \in \varCh' \iff \exists x \in X, x = \{y\}\]
\end{remark}

\noindent
Il existe plusieurs reformulations naturelles de l'axiome du choix. Pour
celles-ci, nous allons introduire le formalisme pour traiter des fonctions.

\begin{definition}[Fonction]
  Soient $X,Y$ deux ensembles. Une fonction partielle
  $\makeFunNamePartial{f}{X}{Y}$ est
  une partie $\varGrP{f}\subseteq X \times Y$ (appelée graphe de $f$) telle que
  le prédicat $(x,y)\in \varGrP{f}$, défini pour $x\in X$ et $y\in Y$, est une
  relation fonctionnelle. On désigne par $\dom(f)$ et $\im(f)$ les ensembles
  associés à la relation fonctionnelle. On dit que $f$ est partielle si
  $\dom(f) \neq X$, et totale si $\dom(f) = X$. Pour une fonction totale
  $f$, on notera $\makeFunName{f}{X}{Y}$, et on appellera simplement
  \og fonction\fg les fonctions totales.

  Pour deux ensembles $X,Y$, on note par $\Funct(X,Y)$ l'ensemble des fonctions
  (totales) de $X$ dans $Y$~:
  \[\Funct(X,Y) \defeq \compr{f \subseteq X \times Y}{\FunTot(``(x,y)\in f")}\]
\end{definition}

\begin{exercise}
  Soient $X$ et $Y$ deux ensembles. Montrer que $\Funct(X,Y)$ est bien un
  ensemble.
\end{exercise}

\begin{remark}
  Dans la présentation faite ici, la notion de totalité est relative à des
  ensembles de départ et d'arrivée fixés. On considère que si $f$ et $g$ ont
  le même graphe, alors $f = g$. Suivant les auteurs, on peut croiser un
  formalisme considérant qu'une fonction $\makeFunName{f}{X}{Y}$ est un triplet
  $(X,Y,\varGrP{f})$. \'Etant donné que nous étudierons presque toujours des
  fonctions dont l'ensemble de départ et d'arrivée sont précisés, ces deux
  présentations sont interchangeables.
\end{remark}

\begin{notation}
  Dans le cas d'une fonction $\makeFunName{f}{X}{Y}$, il est naturel d'ajouter
  un symbole de fonction correspondant, pour externaliser l'objet ensembliste
  interne (et ainsi parler dans notre langage de la fonction $f$ associée).
  Il existe plusieurs façons de traiter formellement cet ajout. On peut
  par exemple ajouter une opération agissant sur tout $\unZF$ de la façon
  suivante~:
  \[\App(X,Y,f,x,y) \defeq x \in X \land y \in Y \land (x,y) \in \varGrP{f}\]
  Cette formule $\App$ définit une fonction partielle qu'on peut rendre
  totale \latinexpr{e.g.} en ajoutant que s'il n'existe aucun $y$ correspondant
  à un $x$ donné, alors $\App(X,Y,f,x,z)$ est vraie pour un $z$
  quelconque (typiquement $\vide$). On peut ainsi ajouter un symbole de fonction
  $\App$ et considérer que $f(x)$ est une réécriture plus lisible du symboles de
  fonction associé à $\App(X,Y,f,x)$.
\end{notation}

\begin{notation}
  Pour définir une fonction, on utilisera généralement la notation
  \[\makeFun{f}{X}{Y}{x}{t}\]
  où $t$ est un terme pouvant faire intervenir la variable $x$, permettant de
  donner une description fiable du graphe $\varGrP{f}$, ainsi que des ensembles
  $X,Y$ respectivement de départ et d'arrivée. Dans le cas où on souhaite
  uniquement désigner le graphe de la fonction, on pourra écrire $x \mapsto t$
  directement.
\end{notation}

\noindent
On définit de plus les notions d'injectivité, de surjectivité et de bijectivité.

\begin{definition}[Injectivité, surjectivité, bijectivité]
  Une fonction $\makeFunName{f}{X}{Y}$ est dite injective si chaque élément
  $y\in Y$ a au plus un antécédent, c'est-à-dire si la proposition suivante est
  vérifiée~:
  \[\forall x_0,x_1\in X, f(x_0) = f(x_1)\implies x_0 = x_1\]

  Une fonction $\makeFunName{f}{X}{Y}$ est dite surjective si tout élément
  $y\in Y$ a au moins un antécédent, c'est-à-dire si la proposition suivante est
  vérifiée~:
  \[\forall y\in Y, \exists x \in X, f(x) = y\]

  Une fonction $\makeFunName{f}{X}{Y}$ est dite bijective si elle est à la fois
  injective et surjective. On notera $\makeFunNameIso{f}{X}{Y}$ lorsqu'on veut
  signifier que la fonction $f$ établit une bijection entre $X$ et $Y$.
\end{definition}

En combinant notre construction de $\bN$ aux considérations fonctionnelles, on
en vient à définir le principe de récursion sur $\bN$.

\begin{exercise}
  Montrer que l'on peut définir des fonctions par récursion sur $\bN$,
  c'est-à-dire montrer que si l'on a un ensemble $X$, un élément $x_0\in X$ et
  une fonction $\makeFunName{f_{\Ss}}{\bN \times X}{X}$ alors il existe une
  unique fonction $\makeFunName{f}{\bN}{X}$ telle que
  \[\begin{cases}
  f(0) = x_0\\
  \forall n \in \bN, f(\Ss n) = f_{\Ss}(n,f(n))
  \end{cases}\]
\end{exercise}

L'exercice suivant permet de définir des fonctions depuis un produit cartésien
en considérant seulement ses projections.

\begin{exercise}
  Soient $X,Y,Z$ des ensembles. Montrer qu'il existe une bijection entre
  $\Funct(X\times Y,Z)$ et $\Funct(X,\Funct(Y,Z))$ donnée par
  $f\mapsto (x\mapsto y \mapsto f(x,y))$.
\end{exercise}

\begin{exercise}
  Montrer que pour tout ensemble $X$, il existe une unique fonction
  $\makeFunName{\mathord{!}}{\vide}{X}$.
\end{exercise}

Donnons maintenant deux nouvelles formulations de l'axiome du choix. Dans ces
formulations, plutôt que d'isoler un ensemble $\varCh$ intersectant chaque
élément d'un ensemble $X$ en des singletons, on considère directement une
fonction qui va, à un ensemble $X$, associer à chaque $x \in X$ un élément de
$x$. On appelle une telle fonction une fonction de choix, et l'axiome du choix
peut se formuler en un principe d'existence de fonctions de choix.

\begin{proposition}
  L'axiome du choix est équivalent à la proposition suivante~:
  \[\forall x,
  \vide\notin x \implies
  \exists \makeFunName{f}{x}{\bigcup x}, \forall y \in x, f(y) \in y\]
\end{proposition}

\begin{proof}
  Supposons l'axiome du choix. On considère un ensemble $x$ dont tous les
  éléments sont non vides. On définit alors l'ensemble
  \[\varCh_0\defeq \compr{\compr{(y,z)}{z \in y}}{y \in x}\]
  Comme $x$ n'a pas d'élément vide, chaque élément de $\varCh_0$ est non vide.
  De plus, si $a \neq b \in \varCh_0$, alors
  $a = \compr{(a_0,y)}{y \in a_0}$ et $b = \compr{(b_0,y)}{y \in b_0}$, donc
  $a\cap b = \vide$ en considérant la première coordonnée de chaque élément.
  On peut donc appliquer l'axiome du choix sur
  $\varCh_0$ pour trouver $\varCh$ tel que pour tout $y\in x$,
  \[\compr{(y,z)}{z \in y}\cap \varCh = \{(y,y_0)\}\]
  donc $\varCh\cap \bigcup \varCh_0$ définit exactement une relation
  fonctionnelle $f$ (à $y$ on n'associe qu'une seule image $y_0$), et par
  définition de $\varCh_0$ on sait que $f(y)\in y$.

  Supposons qu'on possède une fonction de choix pour tout ensemble sans élément
  vide. Soit $X$ un ensemble dont les éléments sont non vides et deux à deux
  disjoints. On peut alors trouver $\makeFunName{f}{X}{\bigcup X}$ telle que
  $f(x)\in x$ pour tout $x\in X$. Soit alors $\varCh = \im(f)$. Comme
  $f(x)\in x$ et tous
  les éléments de $X$ sont deux à deux disjoints, on en déduit que
  $\varCh\cap x$ contient au plus un élément, pour chaque $x\in X$. Comme de
  plus $f(x) \in x$, $\varCh\cap x$ contient au moins un élément. Donc
  $\varCh\cap x$ est un singleton.
\end{proof}

\begin{proposition}
  L'axiome du choix est équivalent à la proposition suivante~:
  \[\forall x, %x\neq\vide\implies
  \exists \makeFunName{f}{\powerset(x) \setminus \{\vide\}}{x},
  \forall y \in \powerset(x)\setminus \{\vide\}, f(y)\in y\]
\end{proposition}

\begin{proof}
  Cette formulation est un cas particulier de la formulation précédente, puisque
  $\powerset(x)\setminus\{\vide\}$ ne contient que des éléments non vides.

  Réciproquement, supposons qu'on a $x$ dont les éléments sont non
  vides. On a alors
  \[x\subseteq \powerset\bigg(\bigcup x\bigg)\setminus\{\vide\}\]
  d'où le résultat en spécialisant une fonction de choix donnée par hypothèse
  sur $x$.
\end{proof}

Une autre formulation importante nécessite la notion de produit, que l'on peut
simplement considérer comme une généralisation de l'opération $\times$.

\begin{definition}[Produit cartésien quelconque]
  Soit $X$ un ensemble, on définit $\prod X$ par
  \[\prod X \defeq
  \compr{\makeFunName{f}{X}{\bigcup X}}{\forall x \in X, f(x)\in x}\]

  Si $\{X_i\}_{i\in I}$ est une famille d'ensembles (c'est-à-dire une fonction
  $\makeFunName{g}{i}{\unZF}$, \textit{i.e.} une relation fonctionnelle entre un
  ensemble $I$ et des ensembles quelconques) alors on considère que
  \[\prod_{i\in I} X_i \defeq
  \compr{\makeFunName{f}{I}
    {\bigcup_{i\in I} X_i}}{\forall i \in I, f(i) \in X_i}\]
\end{definition}

L'axiome du choix peut alors se réécrire en disant que le produit de toute
famille non vide d'ensembles non vides est lui-même non vide.

\begin{proposition}
  L'axiome du choix est équivalent à la proposition suivante :
  \[\forall x, x\neq \vide \land (\forall y \in x, y \neq \vide)
  \implies \prod x \neq \vide\]
\end{proposition}

\begin{proof}
  Par définition, $\prod X$ est non vide si et seulement si on a une fonction
  $\makeFunName{f}{X}{\bigcup X}$ telle que $\forall x \in X, f(x)\in x$. C'est
  un énoncé précédent de l'axiome du choix.
\end{proof}

\subsection{Axiome de fondation}

Le dernier axiome de $\ZFC$ est l'axiome de fondation. Selon les auteurs, il est
ou non ajouté à $\ZFC$~: ici, nous considérerons que $\ZFC$ le contient. Cet
axiome empêche notamment l'existence de cycles d'appartenance, par exemple
d'ensemble $x$ tel que
$x\in x$. Formellement, l'axiome énonce que la relation $\in$ sur l'univers
ensembliste $\unZF$ est bien fondée au sens de la \cref{def.wf.ord}, à ceci
près que la relation n'est pas un ordre.

\begin{axiom}[Fondation]\label{ax.ZF.AF}
  Pour tout ensemble $x$ non vide, il existe $y\in x$ tel que
  $x\cap y = \vide$. Formellement, cela s'écrit :
  \[\forall x, x\neq\vide \implies \exists y \in x, x\cap y = \vide\]
\end{axiom}

\begin{remark}
  La formulation précédente est une façon standard de donner l'axiome de
  fondation. Cependant, celle-ci transmet difficilement l'idée sous-jacente.
  La bonne fondation de $\in$ se traduit, de façon littérale, par le fait que
  pour tout ensemble $x$ il existe un ensemble $y$ minimal pour $\in$ dans $x$.
  Être minimal pour $\in$ dans $x$ signifie qu'il n'existe aucun élément
  $z \in x$ tel que $z \in y$, ce qui revient à dire que $x \cap y = \vide$.
  Ainsi, il existe $y$ tel que $x \cap y = \vide$, mais cette proposition est à
  considérer comme portant avant tout sur la relation $\in$ elle-même.
\end{remark}

Si l'on voit un ensemble $x$ comme un arbre, où chaque n\oe ud est un ensemble
et où l'on place une arête entre $x$ et $y$ si $x\in y$, alors l'axiome de
fondation stipule que pour chaque ensemble, l'arbre ainsi construit n'a pas de
branche infinie (ni de cycle).

Remarquons que l'arbre que nous avons construit contient non seulement les
éléments de $x$ mais aussi les éléments des éléments de $x$, et ainsi de suite.
Cette idée nous permet d'introduire la notion d'ensemble transitif, qui est un
ensemble tel que l'arbre ainsi généré ne contient que des éléments de l'ensemble
(c'est-à-dire que tous les éléments des éléments de $x$ sont aussi éléments de
$x$).

\begin{definition}[Ensemble transitif]
  Soit un ensemble $x$. On dit que $x$ est transitif si tout élément d'un
  élément de $x$ est aussi élément de $x$. Formellement, cela s'écrit
  \[\trans(x) \defeq \forall y\in x, \forall z \in y, z\in x\]
\end{definition}

Remarquons qu'étant donné un ensemble, on peut toujours ajouter ses éléments,
les éléments de ses éléments et ainsi de suite pour obtenir à la fin un ensemble
transitif.

\begin{proposition}
  Soit $x$ un ensemble. Alors il existe un unique plus petit ensemble $\trcl(x)$
  transitif contenant $x$.
\end{proposition}

\begin{proof}
  Pour construire $\trcl(x)$, on construit la suite d'ensembles
  $(x_i)_{i\in\bN}$ suivante :
  \[\begin{cases}
  x_0 \defeq x\\
  x_{i+1} \defeq \bigcup x_i
  \end{cases}\]
  L'ensemble $\trcl(x)$ est alors $\displaystyle\bigcup_{n\in \bN} x_n$.

  Vérifions que $\trcl(x)$ est bien transitif et contient $x$. Comme
  $x_0\subseteq \trcl(x)$ par définition de l'union, $\trcl(x)$ contient $x$.
  Soient $y\in \trcl(x)$ et $z\in y$. Par définition, il existe $n\in \bN$
  tel que $y\in x_n$. Alors $z\in x_{n+1}$ puisque $z\in y \in x_n$ et
  $x_{n+1}=\bigcup x_n$. Ainsi $z\in \trcl(x)$.

  Soit un ensemble transitif $y$ contenant $x$. Alors par récurrence sur $n$,
  $x_n\subseteq y$~:
  \begin{itemize}
  \item par hypothèse, $x\subseteq y$.
  \item si $x_n\subseteq y$, alors soit $z\in x_{n+1}$, par définition on trouve
    $a\in x_n$ tel que $z\in a \in x_n$, et comme $x_n\subseteq y$, $a\in y$.
    Mais $y$ est transitif, donc $z\in y$ : on en déduit que
    $x_{n+1}\subseteq y$.
  \end{itemize}
  Par récurrence, on en déduit que $\forall n\in\bN, x_n\subseteq y$,
  d'où $\trcl(x)\subseteq y$. Comme on a une inclusion, si deux ensembles ont la
  même propriété d'être les plus petits ensembles transitifs contenant $x$,
  alors ils sont inclus l'un dans l'autre et donc égaux, d'où l'unicité.
\end{proof}

\begin{exercise}
  Montrer qu'un ensemble $x$ est transitif si et seulement si la propriété
  suivante est vérifiée~:
  \[\forall y\in x, y\subseteq x\]
\end{exercise}

La clôture transitive est l'outil par excellence permettant de prouver que
certaines équations d'ensembles sont impossibles sous l'axiome de fondation.

\begin{exercise}
  Pour chaque définition de l'ensemble $x$, donner $\trcl(x)$ et montrer que
  l'axiome de fondation empêche l'existence de l'ensemble $x$~:
  \begin{enumerate}[label=(\roman*)]
  \item $x = \{x \}$~;
  \item $x = \{\{\{x \}\}\}$~;
  \item $x = \{y\}$ et $y = \{x\}$.
  \end{enumerate}

  Montrer qu'il n'existe pas de suite décroissante pour l'inclusion,
  c'est-à-dire de famille d'ensembles $(x_i)_{i \in \bN}$ vérifiant
  \[\forall i \in \bN, x_{i+1} \in x_i\]
  \textit{Indication~: on pourra considérer l'ensemble
    $\compr{x_i}{i \in \bN}$.}
\end{exercise}

Les deux propriétés de bonne fondation de $\in$ et de transitivité mènent au
lemme d'effondrement de Mostowski. Celui-ci montre que toute classe munie d'une
relation se comportant suffisamment comme la relation $\in$ peut en fait se
simuler par une unique classe transitive bien fondée où la relation est
directement $\in$.

\begin{lemma}[Effondrement de Mostowski \cite{Mostovski}]\label{lem.most}
  Soit $\varCC$ une classe et $\varRR$ une relation vérifiant :
  \begin{itemize}
  \item $\varRR$ définit des ensembles, c'est-à-dire que pour tout $x\in\varCC$,
    la collection définie par $\LclR{x} \defeq \compr{y}{y \varRR x}$ est
    un ensemble.
  \item $\varRR$ est bien fondée, c'est-à-dire :
    \[\forall X\subseteq \varCC, \exists x \in X, \LclR{x}\cap X = \vide\]
  \item $\varRR$ est extensionnelle, c'est-à-dire :
    \[\forall x,y\in \varRR, \LclR{x} = \LclR{y} \implies x = y\]
  \end{itemize}

  Alors il existe une unique classe $\varCM$ transitive et un unique
  $\makeFunNameIso{(\varCol,\varRR)}{\varCC}{(\varCM,\in)}$, isomorphisme entre
  les deux classes munies d'une relation.
\end{lemma}

\begin{remark}
  Nous n'avons pas formellement défini ce qu'est un isomorphisme entre classes.
  Cette définition est assez naturelle~: étant données deux classes $\varCC$
  et $\varCC'$ avec chacune une relation $\varRR$ (respectivement $\varRR'$), un
  isomorphisme $\makeFunNameIso{\varFunCl}{(\varCC,\varRR)}{(\varCC',\varRR')}$
  est une formule $\varFunCl(x,y)$ à deux variables libres telle que, pour une
  proposition $\varFP$ définissant $\varCC$ et une proposition $\varFC$
  définissant $\varCC'$, les propositions suivantes peuvent être prouvées :
  \begin{itemize}
  \item $\forall x,\forall y,\forall z, \varFP(x)\land \varFC(y)\land\varFC(z)
    \land\varFunCl(x,y)\land\varFunCl(x,z) \implies y = z$
  \item $\forall x,\varFP(x)\implies \exists y, \varFC(y)\land\varFunCl(x,y)$
  \item $\forall x,\forall y,\forall z, \varFP(x)\land\varFP(y)\land\varFC(z)
    \land\varFunCl(x,z)\land\varFunCl(y,z)\implies x = y$
  \item $\forall y,\varFC(y)\implies \exists x, \varFP(x)\land\varFunCl(x,y)$
  \item $\forall x,\forall x',\forall y, \forall y',
    \varFP(x)\land\varFP(x')\land
    \varFC(y)\land\varFC(y')\land\varFunCl(x,y)\land\varFunCl(x',y')\implies
    (x\varRR x'\iff y\varRR' y')$
  \end{itemize}

  Cependant, nous allons simplement travailler sur les classes comme sur les
  ensembles, puisque les manipulations syntaxiques correspondant à nos
  arguments ensemblistes habituels sont les mêmes que celles que nous allons
  faire sur nos classes. Remarquons simplement que ce lemme est, puisqu'il parle
  de classe, un schéma de lemmes paramétré par les propositions définissant
  $\varCC$ et $\varRR$.
\end{remark}

\begin{proof}
  Montrons d'abord que si $\varCC_1$ et $\varCC_2$ sont deux classes transitives
  et que $\makeFunNameIso{\varFunCl}{\varCC_1}{\varCC_2}$ est un isomorphisme
  pour la relation $\in$, alors $\varFunCl$ est l'identité.

  Par l'absurde, supposons que $\varFunCl$ n'est pas l'identité. On pose
  $\varCC_0 \defeq \compr{x}{\varFunCl(x)\neq x}$, non vide par hypothèse. Soit
  alors $x_0 \in \varCC_0$~: $\trcl(\{x_0\}) \cap \varCC_0$ est un ensemble non
  vide (il contient au moins $x_0$) et, par
  fondation, il est possible de choisir $x$ minimal, \latinexpr{i.e.}
  $x \in \trcl(\{x_0\})\cap \varCC_0$ tel que
  $x\cap \trcl(\{x_0\})\cap \varCC_0 = \vide$. Alors,
  pour tout $y \in x$, on sait (par minimalité de $x$) que $y\notin \varCC_0$,
  donc par définition que $\varFunCl(y) = y$. Mais alors, comme $\varFunCl$ est
  un isomorphisme pour $\in$~:
  \begin{align*}
    \varFunCl(x) &= \{z \in \varFunCl(x)\} \\
    &= \compr{\varFunCl(y)}{\varFunCl(y) \in \varFunCl(x)} &
    \text{en posant}\; y \defeq \theta^{-1}(z) \\
    &= \compr{\varFunCl(y)}{y \in x} \\
    &= \compr{y}{y \in x} & \text{par la remarque précédente}\\
    \varFunCl(x) &= x & \text{par extensionalité}
  \end{align*}
  donc $\varFunCl(x) = x$, ce qui contredit le fait que $x$ soit dans $C$.

  Reprenons maintenant notre classe $\varCC$. On définit alors la classe
  $\varCM$ par l'image de la fonction $\varCol$ définie comme suit :
  \[
  \makeFun{\varCol}{\varCC}{\unZF}{x}{\compr{\varCol(y)}{y\in\LclR{x}}}
  \]
  c'est-à-dire que la relation $\varCM$ est définie par la proposition
  \[\varCM(x,z) \defeq z = \compr{\varCol(y)}{y \in \LclR{x}}\]
  où l'on prouve par l'absurde que cette proposition vérifie les propriétés d'un
  isomorphisme entre $(\varCC,\varRR)$ et $(\varCM,\in)$, en utilisant un
  principe analogue à la preuve précédente. La classe $\varCM$ est définie
  par la proposition
  \[z \in \varCM \defeq \exists x\in\varCC, \varCol(x,z)\]

  L'unicité découle alors du fait que s'il existe une autre telle classe
  $\varCM'$, alors l'isomorphisme entre elle et $\varCM$ vérifie les
  hypothèses du résultat intermédiaire, et est donc l'identité. D'où
  $\varCM = \varCM'$.
\end{proof}

\begin{remark}
  Notre définition de $\varCol$ peut sembler illégitime, étant donné que le
  \cref{thm.recur.bf} ne couvre que le cas où le domaine est un ensemble (et
  qu'il est un ordre bien fondé, mais il est assez clair que cela se déroule de
  la même façon pour une relation bien fondée qui n'est pas forcément un ordre).
  Cependant, on voit que pour définir $\varCol(y)$, il suffit de le définir sur
  un ensemble, celui des éléments inférieurs (comme $\varRR$ définit des
  ensembles, cela fonctionne bien).

  Ainsi la version purement formelle de $\varCol$ devrait être, avec $\varFF$ le
  prédicat définissant $\varCC$~:
  \begin{multline*}
    \varCol(x,y) \defeq \exists X,
    (\forall a \in X, \varFF(a) \land \forall b, b\varRR a \implies b \in X)
    \land \exists \makeFunName{f}{X}{\unZF}, \\(\forall z\in X,
    \forall Y, (\forall a, a \in Y \iff a \varRR z) \implies
    f(z) = \compr{f(a)}{a \varRR z})\land f(x) = y
  \end{multline*}
  et l'utilisation du \cref{thm.recur.bf} permet de conclure que cette relation
  est bien fonctionnelle.

  Au sein même de cette définition, on voit apparaître l'écriture
  $\makeFunName{f}{X}{\unZF}$, qui peut là encore sembler litigieuse. Une
  fonction $X \to \unZF$ est en fait un ensemble décrivant une relation
  fonctionnelle totale dont le domaine est $X$, sans restriction sur l'ensemble
  d'arrivée. Ces fonctions $X \to \unZF$ constituent une classe, mais
  cela n'est pas un problème lorsqu'on effectue une quantification.
\end{remark}

\begin{remark}
  Dans la conclusion de la démonstration, dire que $\varCM = \varCM'$
  signifie que $x \in \varCM \iff x \in \varCM'$, \latinexpr{i.e.} que
  les formules $\varFF$ et $\varFF'$ définissant respectivement $\varCM$
  et $\varCM'$ sont équivalentes.
\end{remark}

\subsection{Résumé~: les fragments de ZFC}\label{subsec.frag.ZFC}

Concluons cette section en redonnant les différents axiomes de $\ZFC$. En fait,
en enlevant les doublons, il nous suffit d'un petit fragment de ce que nous
avons vu. Donnons d'abord les axiomes de $\ZF$~:
\begin{itemize}
\item l'\cref{ax.ZF.ext} (extensionalité) (noté $\AxExt$)
\item l'\cref{ax.ZF.union} (réunion) (noté $\AxUn$)
\item l'\cref{ax.ZF.pow} (ensemble des parties) (noté $\AxPow$)
\item l'\cref{ax.ZF.repl} (schéma de remplacement) (noté $\AxRepl$)
\item l'\cref{ax.ZF.infini} (infini) (noté $\AxInf$)
\item l'\cref{ax.ZF.AF} (fondation) (noté $\AxF$)
\end{itemize}
$\ZFC$ est alors obtenue comme $\ZF + \AxC$. On peut aussi définir
$\ZF - \AxF$ par exemple. De plus, la théorie $\Zermelo$, introduite par
Ernst \name{Zermelo} dans \cite{zermelo08}, est la
théorie contenant le schéma d'axiomes de compréhension mais pas de remplacement,
ne contenant pas l'axiome de fondation et contenant l'axiome de la paire (et les
axiomes d'extensionalité, de réunion, d'ensemble des parties et de l'infini).
Dans la suite du chapitre, nous n'aurons besoin que de $\Zermelo - \AxF$.

\section{Construction des ensembles de nombres}

Nous avons construit $\bN$ directement grâce à l'axiome de l'infini, mais
celui-ci n'est pas le seul ensemble de nombres important à construire. Dans
cette section, nous allons voir la construction de $\bZ, \bQ$ et
$\bR$, ainsi que leurs propriétés essentielles. Nous allons commencer
notre étude par celle de $\bN$, puisque nous avons simplement défini cet
ensemble et le principe de récurrence.

\subsection{Les entiers naturels}

Rappelons les points essentiels à propos de $\bN$~: cet ensemble contient
$0$, une fonction $\Ss$ unaire, et le principe de récurrence. De plus, on peut
définir une fonction par récursion. Nous pouvons donc définir les fonctions
$+$ et $\times$.

\begin{definition}[Addition, multiplication]
  On définit la fonction $\makeFunName{+}{\bN \times \bN}{\bN}$ et la
  fonction $\makeFunName{\times}{\bN \times \bN}{\bN}$ par récursion
  sur leur second argument~:
  \begin{itemize}
  \item $\forall n \in \mathbb N, n + 0 = n$
  \item $\forall n,m\in \mathbb N, n + (\Ss m) = \Ss(n + m)$
  \item $\forall n \in \mathbb N, n \times 0 = 0$
  \item $\forall n,m\in\mathbb N, n \times (\Ss m) = n\times m + n$
  \end{itemize}
\end{definition}

\begin{notation}
  On définit $1 = \Ss 0$, $2 = \Ss\Ss 0$ et ainsi de suite. Plutôt que $\Ss n$,
  nous utiliserons souvent l'expression $n + 1$.
\end{notation}

\begin{property}
  Les propriétés suivantes sont vérifiées~:
  \begin{itemize}
  \item $0$ est un élément neutre pour $+$~: pour tout $n\in \bN$,
    $0 + n = n$ et $n + 0 = n$.
  \item $+$ est commutatif~: pour tous $n,m\in\bN$, $n+m=m+n$.
  \item $+$ est associatif~: pour tous $n,m,p\in\bN$,
    $n+(m+p) = (n+m)+p$.
  \item $+$ est régulier~: pour tous $n,p,m\in\bN$, si $n+m = n+p$ alors
    $m=p$.
  \item pour tous $n,m\in \bN$, si $n + m = 0$ alors $n=0$ et $m = 0$.
  \item $+$ distribue sur $\times$~: pour tous $n,m,p\in\bN$,
    $n\times(m+p) = n\times m + n \times p$ et
    $(n+m)\times p = n\times p + m \times p$.
  \item $1$ est un élément neutre pour $\times$.
  \item $0$ est absorbant pour $\times$~: pour tout $n\in \bN$,
    $0\times n = 0$.
  \item $\times$ est commutatif.
  \item $\times$ est associatif.
  \item La structure est intègre~: pour tous $n,m\in \bN$, si
    $n\times m = 0$ alors $n = 0$ ou $m = 0$.
  \end{itemize}
\end{property}

\begin{proof}
  On montre seulement quelques résultats parmi ceux-ci~:
  \begin{itemize}
  \item $0$ est neutre~: pour la première égalité, le résultat est par
    définition. Pour la deuxième, montrons par récurrence sur $n$ que
    $0+n = n$~:
    \begin{itemize}
    \item $0 + 0 = 0$ par définition de $+$.
    \item supposons que $0+n = n$, alors $0+\Ss n = \Ss(0+n) = \Ss n$.
    \end{itemize}
    D'où le résultat par récurrence.
  \item $+$ est commutatif~: on montre d'abord par récurrence sur $m$ que pour
    tous $n,m\in\bN$, on a $\Ss n+m = \Ss (n+m)$~:
    \begin{itemize}
    \item $\Ss n+0=\Ss n = \Ss(n+0)$.
    \item si $\Ss (n+m) = \Ss(n+m)$, alors
      \begin{align*}
        \Ss n+\Ss m &= \Ss(\Ss n+m)\\
        &= \Ss(\Ss(n+m))\\
        &= \Ss(n+\Ss m)
      \end{align*}
    \end{itemize}
    Nous donnant un premier résultat par récurrence. On montre maintenant que
    $+$ est commutatif par récurrence sur l'argument de droite~:
    \begin{itemize}
    \item $n+0 = n = 0+n$ comme nous l'avons déjà prouvé.
    \item supposons que $n+m = m+n$, alors
      \begin{align*}
        n + S\;m &= S\;(n+m)\\
        &= S\;(m+n)\\
        &= S\;m + n
      \end{align*}
      en utilisant le résultat précédent.
    \end{itemize}
    Donc $+$ est commutatif.
  \item $+$ est associatif. On prouve par récurrence sur $k$ la proposition
    \[\forall k\in\bN,\forall n\in\bN,\forall m\in\bN,n+(m+k)=(n+m)+k\]
    \begin{itemize}
    \item si $k = 0$ alors $n + (m + 0) = n + m$ et $(n+m)+0=n+m$, d'où
      l'égalité.
    \item supposons la proposition pour $k$, alors
      \begin{align*}
        n+(m+\Ss k) &= n + \Ss (m + k)\\
        &= \Ss(n+(m+k))\\
        &= \Ss((n+m)+k)\\
        &= (n+m)+\Ss k
      \end{align*}
    \end{itemize}
    D'où le résultat par récurrence.
  \end{itemize}
  L'ensemble des propriétés confère à $(\bN,0,1,+,\times)$ une structure
  de demi-anneau intègre.
\end{proof}

\begin{exercise}
  Démontrer les autres propriétés non démontrées ci-dessus.
\end{exercise}

\noindent
La relation $\leq$ est aussi définissable grâce à notre langage ensembliste.

\begin{definition}[Inégalité]
  On définit la relation $\mathord{\leq}\subseteq \bN \times \bN$
  par~:
  \[n \leq m \defeq \exists k \in \bN, m = n + k\]
\end{definition}

\begin{property}
  La relation $\leq$ est une relation d'ordre.
\end{property}

\begin{proof}
  Prouvons que la relation est réflexive, antisymétrique et transitive~:
  \begin{itemize}
  \item soit $n\in \bN$, alors $n = n + 0$, d'où $n\leq n$.
  \item soient $n,m\in \bN$, supposons que $n \leq m$ et $m \leq n$. On
    trouve alors $k,k'$ tels que $n = m + k$ et $m = n + k'$. En substituant
    la deuxième inégalté dans la première, on obtient alors
    \begin{align*}
      n &= (n + k') + k\\
      &= n + (k' + k)\\
      n + 0 &= n + (k' + k)\\
      0 &= k + k'&\text{par régularité de }+\\
    \end{align*}
    d'où $0 = k$, donc $n = m$.
  \item soient $n,m,p\in \bN$ tels que $n\leq m$ et $m\leq p$. On trouve
    donc $k,k'$ tels que $m = n + k$ et $p = m + k'$. Par substitution, on en
    déduit que $p = (n + k) + k'$, et par associativité $p = n + (k + k')$, d'où
    $n \leq p$.
  \end{itemize}
  Ainsi, $\leq$ est une relation d'ordre.
\end{proof}

\subsection{Les entiers relatifs}

Nous allons maintenant construire les entiers relatifs, dont l'ensemble est
noté $\bZ$, avec la construction classique~: on cherche à transformer
le monoïde $(\bN,0,+)$ en un groupe. Cette construction se généralise
en fait en une construction sur n'importe quel monoïde commutatif.

Nous allons donc voir d'abord la construction de $\bZ$, puis la
construction plus générale du symétrisé d'un monoïde commutatif.

L'ensemble $\bZ$ sera vu comme un quotient de l'ensemble
$\bN \times \bN$. L'idée est de considérer une paire $(n,m)$ comme
l'entier $n - m$. On voit alors que pour avoir $(n,m) = (p,q)$, il faut et il
suffit que $n - m = p - q$, c'est-à-dire que $n + q = p + m$. Cela motive la
définition de notre relation $\sim$.

\begin{definition}[Congruence additive]
  On définit la relation binaire $\sim$ sur $\bN \times \bN$ par~:
  \[(n,m)\sim (p,q) \defeq n + q = p + m\]
\end{definition}

\begin{proposition}
  La relation $\sim$ est une relation d'équivalence.
\end{proposition}

\begin{proof}
  On vérifie les axiomes d'une relation d'équivalence~:
  \begin{itemize}
  \item on voit que $n+m=n+m$ donc $(n,m)\sim (n,m)$.
  \item supposons que $(n,m) \sim (p,q)$, alors $n + q = p + m$, donc par
    symétrie de l'égalité $p + m = n + q$, donc $(p,q)\sim (n,m)$.
  \item supposons que $(n,m)\sim(p,q)$ et $(p,q)\sim (r,s)$. On sait donc que
    $n + q = p + m$ et $p + s = r + q$. En ajoutant $r$ de chaque côté dans
    la première équation et en réordonnant les termes, on a alors
    \[n + (r + q) = p + m + r\]
    d'où $n + p + s = p + m + r$ donc, par régularité, $n + s = r + m$.
    Donc $(n,m)\sim(r,s)$.
  \end{itemize}
  Ainsi $\sim$ est une relation d'équivalence.
\end{proof}

Pour simplifier nos constructions, nous allons maintenant introduire la notion
de congruence.

\begin{definition}[Congruence]
  Soit $(X,\star)$ un magma. On appelle congruence une relation d'équivalence
  $\mathord{\sim}\subseteq X \times X$ telle que
  \[\forall x,y,x',y' \in X, x\sim x' \land y \sim y'
  \implies x\star y \sim x' \star y'\]
\end{definition}

Une congruence nous permet de définir directement une opération sur l'ensemble
quotient.

\begin{proposition}
  Soit $(X,\star)$ un magma et $\sim$ une congruence sur $(X,\star)$. Alors
  il existe une unique opération de magma $\star'$ définie sur $\quot{X}{\sim}$
  vérifiant
  \[\forall x,y \in X, \clas{x\star y} = \clas{x}\star' \clas{y}\]
\end{proposition}

\begin{proof}
  Le fait qu'il existe au plus une telle opération découle directement du fait
  que la propriété définit l'opération $\star'$. On veut donc montrer que
  l'opération est bien définie~: soient $x,x',y,y'\in X$ tels que
  $x\sim x'$ et $y\sim y'$, alors comme $\sim$ est une congruence, on en
  déduit que $x\star y \sim x' \star y'$, donc que
  \[\clas{x} \star' \clas{y} = \clas{x'}\star' \clas{y'}\]
  c'est-à-dire que le résultat de $\star'$ ne dépend pas du représentant choisi.
\end{proof}

\begin{exercise}
  Montrer que si les propriétés suivantes sont vraies pour $(X,\star)$ alors
  elles le sont aussi pour $(\quot{X}{\sim},\star')$~:
  \begin{itemize}
  \item associativité
  \item commutativité
  \item existence d'un élément neutre
  \item existence d'un inverse
  \item régularité à gauche (respectivement à droite)
  \end{itemize}
\end{exercise}

En revenant sur le cas de notre construction de $\bZ$, on voit que
$\bN\times\bN$ peut être muni d'une structure de monoïde en
considérant l'addition coordonnée par coordonnée~: $(n,m)+(p,q) = (n+p,m+q)$.
Il se trouve que $\sim$ est une congruence sur cette opération.

\begin{proposition}
  $\sim$ est une congruence sur $\bN \times \bN$ avec la structure
  de monoïde coordonnée par coordonnée.
\end{proposition}

\begin{proof}
  Il s'agit d'une simple vérification~: supposons que $(n,m)\sim(n',m')$ et que
  $(p,q)\sim(p',q')$, alors $(n+p,m+q)\sim(n'+p',m'+q')$~:
  \begin{align*}
    n + p + m' + q' &= (n + m') + (p + q')\\
    &= (n' + m) + (p' + q)\\
    &= n' + p' + m + q
  \end{align*}
\end{proof}

Ainsi, en notant $\bZ = \quot{(\bN \times \bN)}{\sim}$, on a
défini une opération $+$ lui conférant une structure de monoïde commutatif. Il
nous reste à prouver que c'est en fait un groupe abélien.

\begin{proposition}
  $(\bZ,+)$ est un groupe abélien.
\end{proposition}

\begin{proof}
  Soit $x \in \bZ$, par définition on trouve $n,m\in\bN$ tels que
  $x = \clas{(n,m)}$, et on pose $-x = \clas{(m,n)}$. On montre que
  $x + (-x) = 0_{\bZ}$. Par définition,
  \[x + (-x) = \clas{(n + m, m + n)}\]
  et $n + m + 0 = 0 + m + n$, donc $\clas{(n+m,m+n)}=\clas{(0,0)}$,
  d'où le résultat.
\end{proof}

\noindent
On s'attarde maintenant à la structure multiplicative de $\bZ$.

\begin{definition}
  On définit sur $\bN\times \bN$ l'opération $\times$ par
  \[(n,m)\times(p,q) = (n\times p + m \times q, m \times p + n \times q)\]
  et $\times$ est définie sur $\bZ$ par passage au quotient de cette
  opération par $\sim$.
\end{definition}

\begin{exercise}
  Montrer que $\sim$ est bien une congruence pour $\times$.
\end{exercise}

\begin{exercise}
  Montrer que $(\bZ,0,+,1,\times)$ est un anneau.
\end{exercise}

On montre qu'on peut injecter le demi-anneau $(\bN,0,+,1,\times)$ dans
l'anneau $(\bZ,0,+,1,\times)$.

\begin{proposition}
  La fonction
  \[
  \makeFun{\iota}{\bN}{\bZ}{n}{\clas{(n,0)}}
  \]
  est un morphisme de demi-anneau injectif.
\end{proposition}

\begin{proof}
  On vérifie chaque axiome~:
  \begin{itemize}
  \item pour tous $n,m\in\bN$~:
    \begin{align*}
      \iota(n) + \iota(m) &= \clas{(n,0)} + \clas{(m,0)}\\
      &= \clas{(n+m,0+0)}\\
      &= \iota(n+m)
    \end{align*}
  \item pour tous $n,m\in\bN$~:
    \begin{align*}
      \iota(n)\times \iota(m) &= \clas{(n,0)} \times \clas{(m,0)}\\
      &= \clas{(n \times m + 0, 0 + 0)}\\
      &= \clas{(n\times m,0)}\\
      &= \iota(n+m)
    \end{align*}
  \item le neutre de l'anneau $\bZ$ est $\clas{(1,0)}$, qui est bien
    $\iota(1)$.
  \item si $\iota(n) = \iota(m)$, alors $\clas{(n,0)} = \clas{(m,0)}$,
    donc $(n,0)\sim (m,0)$, c'est-à-dire $n = m$.\qedhere
  \end{itemize}
\end{proof}

Enfin, sur $\bZ$, on peut caractériser les éléments comme des entiers
naturels ou des opposés d'entiers naturels.

\begin{proposition}
  Pour tout $k \in \bZ$, il existe $n \in \bN$ tel que soit
  $k = \iota(n)$, soit $k = (-\iota(1)) \times \iota(n)$, que l'on écrira
  plus rapidement $k = n$ (respectivement $k = -n$).
\end{proposition}

\begin{proof}
  Soit $k \in \bZ$. Par définition, on trouve $n,m\in\bN$ tels que
  $k = \clas{(n,m)}$. Comme $\leq$ est total, on sait que soit
  $n \leq m$, soit $m \leq n$. Deux cas se présentent alors~:
  \begin{itemize}
  \item si $n \leq m$, alors $m = (m - n) + n$, donc $k = \clas{(0,m - n)}$,
    nous donnant $k = (-\iota(1))\times \iota(m-n)$.
  \item si $m \leq n$, alors $n = (n - m) + m$, donc
    $k = \clas{(n - m, 0)}$, nous donnant $k = \iota(n-m)$.\qedhere
  \end{itemize}
\end{proof}

\noindent
Passons maintenant à la construction générale sur un monoïde commutatif.

\begin{definition}
  Soit $(\varMMo,\opM,\neuMM)$ un monoïde commutatif. On définit sur
  $\varMMo\times\varMMo$ muni des opérations coordonnée par coordonnée la
  congruence suivante~:
  \[(m,m') \sim (a,a') \defeq \exists k \in \varMMo,
  m \opM a' \opM k = m' \opM a \opM k\]

  On appelle symétrisé de $(\varMMo,\opM,\neuMM)$ le groupe abélien obtenu à
  partir de $\varMMo\times \varMMo$ par quotient par $\sim$.
\end{definition}

\begin{proof}
  Montrons que $\sim$ est une congruence~:
  \begin{itemize}
  \item $\sim$ est réflexive~: pour $(m,m')$, on a
    $m \opM m' \star \neuMM = m \opM m' \opM \neuMM$.
  \item $\sim$ est symétrique par symétrie de l'égalité.
  \item $\sim$ est transitive~: supposons qu'on a $(m,m')\sim (a,a')\sim(b,b')$
    et montrons qu'alors $(m,m')\sim (b,b')$. D'après les définitions, on trouve
    $k,k'\in \varMMo$ tels que
    \[\begin{cases}
    m \opM a' \opM k = m' \opM a \opM k\\
    a \opM b' \opM k' = a' \opM b \opM k'
    \end{cases}\]
    En \og additionnant\fg les deux lignes, on obtient
    \[m\opM(a\opM a'\opM k\opM k')\opM b'=m'\opM(a\opM a'\opM k\opM k')\opM b\]
    d'où le résultat en prenant $k'' = a \opM a' \opM k \opM k'$.
  \item $\sim$ est une congruence~: supposons que $(m,m')\sim(a,a')$ et
    $(n,n')\sim(b,b')$, alors
    $(m\opM n, m' \opM n') \sim(a\opM b,a'\opM b')$. En effet, si on
    considère $k,k'\in \varMMo$ tels que
    \[\begin{cases}
    m \opM a' \opM k = m' \opM a \opM k\\
    n \opM b' \opM k' = n' \opM b \opM k'
    \end{cases}\]
    alors on obtient en \og additionnant\fg les deux lignes
    \[(m \opM n) \opM (a' \opM b') \opM (k \opM k') =
    (m' \opM n') \opM (a \opM b) \opM (k \opM k')\]
    nous donnant le résultat en prenant $k'' = k \opM k'$.
  \end{itemize}
  On en déduit que $\sim$ est une congruence. Ainsi le symétrisé est un monoïde
  commutatif.

  Il nous reste à prouver que c'est un groupe. Soit $\clas{(m,m')}$ un
  élément du symétrisé. On définit son inverse par $\clas{(m',m)}$~:
  \[\clas{(m,m')}\opM \clas{(m',m)} =
  \clas{(m\opM m',m'\opM m)}\]
  or $(m,m)\sim (e,e)$ assez directement, d'où le fait que leur produit vaut
  $\clas{(e,e)}$, le neutre du symétrisé.
\end{proof}

\begin{remark}
  On ajoute dans notre définition de $\sim$ un terme $k$~: celui-ci est
  nécessaire pour le cas où notre monoïde n'est pas régulier. On voit en effet
  que ce $k$ absorbe le terme $a \opM a'$ pour la transitivité.
\end{remark}

\begin{exercise}
  Soit $\varMMo$ un monoïde commutatif et $\varMMo^\dagger$ son symétrisé.
  Montrer que la fonction
  \[
  \makeFun{\iota}{\varMMo}{\varMMo^\dagger}{m}{\clas{(m,e)}}
  \]
  est un morphisme de monoïde.
\end{exercise}

\begin{exercise}
  Soit $\varMMo$ un monoïde commutatif et $\varMMo^\dagger$ son symétrisé.
  Montrer que pour tout groupe $\varGrp$ et tout morphisme de monoïde
  $\makeFunName{f}{\varMMo}{\varGrp}$ il existe un unique
  morphisme de groupe $\makeFunName{\relev{f}}{\varMMo^\dagger}{\varGrp}$ tel que
  $f = \relev{f} \circ \iota$, ce
  que l'on peut représenter par le fait que tout diagramme comme suit peut être
  complété en un diagramme commutatif~:
  \begin{center}
    \begin{tikzcd}
      \varMMo \ar[r,"f"]\ar[d,"\iota"] & \varGrp\\
      \varMMo^\dagger \ar[ur,"\hat f",dashed]
    \end{tikzcd}
  \end{center}
\end{exercise}

\subsection{Les nombres rationnels}

Nous avons vu comment construire les nombres entiers en tant que cas particulier
de symétrisation~: $\bZ$ est $\bN$ auquel on a ajouté des inverses pour la loi
$+$. Pour le cas de $\bQ$, on va montrer qu'on peut construire cet ensemble
comme cas particulier du procédé dit de localisation, qui permet d'inverser pour
la loi $\times$ dans un anneau. Comme précédemment, on commence par construire
le cas particulier de $\bQ$ et sa structure de corps avant d'étudier le cas plus
général de la localisation.

Si la construction de $\bZ$ peut paraître assez difficile à visualiser au début,
celle de $\bQ$ sera au contraire très naturelle~: elle consiste simplement à
étudier les fractions. On sait en effet manipuler les fractions comme des
\og couples\fg de la forme $\dfrac{p}{q}$ où plusieurs règles nous permettent de
transformer un couple en un autre. Ces règles peuvent se synthétiser en
\[\frac{p}{q} = \frac{p'}{q'} \defeq p\times q' = p' \times q\]

\`A partir de là, le procédé est très similaire au précédent~: on considère cela
comme une relation d'équivalence sur des couples, et on en déduit la
construction du corps des fractions.

\begin{definition}
  On définit la relation d'équivalence $\sim$ sur
  $\bZ \times (\bN\setminus\{0\})$ par~:
  \[(p,q) \sim (p',q') \defeq p\times q' = p'\times q\]
\end{definition}

\begin{proof}
  On vérifie que cela définit bien une relation d'équivalence~:
  \begin{itemize}
  \item pour tout $(p,q)$, $p\times q = p \times q$ donc $(p,q)\sim(p,q)$.
  \item pour tous $(p,q)\sim (p',q')$, la symétrie de l'égalité nous fait
    déduire que $(p',q')\sim(p,q)$.
  \item si $(p,q)\sim(p',q')\sim(p'',q'')$, alors on sait que
    \[\begin{cases}
    p\times q' = p' \times q\\
    p' \times q'' = p'' \times q'
    \end{cases}\]
    si l'un des nombres parmi $p,p',p''$ est nul, alors tous les autres le sont
    (puisque $q,q',q''$ sont tous les trois non nuls). On en déduit donc qu'on
    peut simplifier (puisque $\bZ$ est un anneau intègre) le produit des
    deux lignes par $q'\times p'$, nous donnant~:
    \[p\times q'' = p'' \times q\]
    donc $(p,q)\sim(p'',q'')$.
  \end{itemize}
  Donc $\sim$ est une relation d'équivalence.
\end{proof}

\begin{exercise}\label{exo.ax.Qring}
  Montrer qu'on peut munir $\bZ\times (\bN\setminus\{0\})$ d'une
  structure d'anneau en définissant les opérations
  \[(p,q)+(p',q') \defeq (p\times q' + p' \times q, q\times q')\qquad
  (p,q)\times(p',q')\defeq (p\times p',q\times q')\]
\end{exercise}

\noindent
Cette structure passe au quotient par $\sim$, ce qui nous donnera la structure
souhaitée pour $\bQ$.

\begin{proposition}
  La relation $\sim$ est une congruence pour les deux opérations définies dans
  l'\cref{exo.ax.Qring}.
\end{proposition}

\begin{proof}
  Soient $(p,q)\sim(p',q')$ et $(r,s)\sim(r',s')$, montrons les deux relations
  \[(p\times s + q\times r, q \times s) \sim
  (p'\times s' + q'\times r',q'\times s')\qquad
  (p\times r, q\times s)\sim(p'\times r', q'\times s')\]
  la seconde relation est la plus facile à montrer~: cela revient à l'égalité
  \[p\times r \times q' \times s' = q\times s \times p' \times r'\]
  mais on sait que $p\times q' = p' \times q$ et que
  $r \times s' = t' \times s$, ce qui nous donne directement le résultat en
  multipliant les deux égalités entre elles.

  Pour la première relation, celle-ci revient à montrer l'égalité
  \[(p\times s + q \times r)\times q' \times s' =
  (p'\times s' + q' \times r')\times q \times s\]
  ce qui, par distributivité, donne
  \[p\times q' \times s \times s' + r \times q \times q' \times s' =
  p'\times q \times s \times s' + r' \times q \times q' \times s\]
  or $p\times q' = p' \times q$ et $r\times s' = r' \times s$, donnant le
  résultat en multipliant les équations respectivement par $s\times s'$ et par
  $q\times q'$, puis en les additionnant.
\end{proof}

\begin{notation}
  On notera $p/q$ ou $\dfrac{p}{q}$ pour noter $\clas{(p,q)}$.
\end{notation}

On peut donc définir $\bQ$~: c'est l'anneau quotient obtenu à partir de $\sim$.
On montre maintenant que cet anneau est un corps.

\begin{proposition}
  $(\bQ,0,+,1,\times)$ est un corps.
\end{proposition}

\begin{proof}
  On a déjà prouvé qu'on avait un anneau commutatif, il nous reste à montrer que
  c'est un corps. Pour cela, soit $p/q\in\bQ$ non nul. Cela signifie en
  particulier que $p \neq 0$. Si $p > 0$, alors on définit l'inverse
  de $p/q$ par $q/p$, et on voit que
  \[\frac{p}{q}\times\frac{p}{q} = \frac{pq}{pq} = 1\]
  Si $p < 0$, alors on définit l'inverse de $p/q$ par $(-q)/(-p)$, et on
  trouve alors
  \[\frac{p}{q}\times\frac{-q}{-p} = \frac{-pq}{-pq} = 1\]

  Dans les deux cas, si $r\in \bQ$ est non nul, alors il est inversible
  pour $\times$.
\end{proof}

\begin{exercise}
  Montrer que la fonction
  \[\makeFun{\iota}{\bZ}{\bQ}{k}{\frac{k}{1}}\]
  est un morphisme d'anneau injectif.
\end{exercise}

Ainsi $\bQ$ peut se voir comme le plus petit corps contenant $\bZ$,
voire comme le plus petit corps contenant $\bN$.

On passe maintenant au cas général. Tout d'abord, on peut voir que la
construction ne demande, pour être effectuée, que d'avoir une partie stable
par $\times$ dans $\bZ$. Si l'on imagine prendre l'ensemble
$\compr{2^n}{n \in \bN}$, par exemple, on peut vérifier que la
construction précédente donne encore une structure, qui ne sera cependant pas
un corps~: ce sera un anneau contenant $\bZ$ dans lequel chaque élément
de $\compr{2^n}{n \in \bN}$ devient inversible. C'est ce procédé qu'on
appelle localisation.

\begin{definition}
  Soit $\varAA$ un anneau commutatif intègre. On considère une partie
  $S\subseteq \varAA$ close par produit, c'est-à-dire telle que
  \[\forall x,y\in S, x \times y \in S\]
  On définit la relation $\sim_S$ sur $\varAA\times S$ par
  \[(a,s) \sim (a',s') \defeq a\times s' = a' \times s\]
\end{definition}

\begin{proposition}
  $\sim_S$ est une congruence pour les opérations
  \[(a,s)+(a',s') \defeq (a\times s' + a' \times s, s\times s')\qquad
  (a,s)\times (a',s') \defeq (a\times a',s\times s')\]
\end{proposition}

\begin{proof}
  On montre les différents axiomes~:
  \begin{itemize}
  \item pour tout $(a,s)$, on a $a\times s = a \times s$ donc
    $(a,s)\sim_S(a,s)$.
  \item la symétrie vient de celle de l'égalité.
  \item si $(a,s)\sim_S(a',s')\sim_S(a'',s'')$, alors $(a,s)\sim_S(a'',s'')$. En
    effet, l'hypothèse nous dit que
    \[\begin{cases}
    a \times s' = a' \times s\\
    a' \times s'' = a'' \times s'
    \end{cases}\]
    d'où, en multipliant les deux lignes~:
    \[a\times (s'\times a') \times s'' = a''\times (s'\times a') \times s\]
    si l'un parmi $a,a',a''$ est nul, les trois sont nuls et le résultat est
    direct. Sinon, par intégrité, on peut simplifier la multiplication par
    $s'\times a'$, nous donnant $a\times s'' = a'' \times s$, donc
    $(a,s) \sim_S (a'',s'')$.
  \item si $(a,s)\sim_S(a',s')$ et $(b,t)\sim_S(b',t')$, pour montrer que
    $(a,s)\times(b,t)\sim_S (a',s')\times (b',t')$, il nous faut montrer
    l'égalité
    \[(a\times t + b \times s)\times s' \times t' =
    (a'\times t' + b' \times s')\times s \times t\]
    et le raisonnement donné plus tôt dans le cas de $\bQ$ fonctionne
    encore.
  \item si $(a,s)\sim_S(a',s')$ et $(b,t)\sim_S(b',t')$, on veut montrer que
    \[a\times b \times s' \times t' = a' \times b' \times s \times t\]
    ce qui se déduit directement des hypothèses.
  \end{itemize}
  Ainsi $\sim_S$ est bien une congruence pour les opérations considérées.
\end{proof}

On peut alors montrer que $\quot{\varAA}{\sim_S}$, qu'on note en général
$S^{-1}\varAA$, contient $\varAA$ et des inverses pour chaque élément de $S$.
En réalité, on peut construire la localisation sans avoir l'intégrité, en
prenant une définition analogue à celle où on ajoute $k$ dans la définition de
la symétrisation.

\subsection{Les nombres réels}

Nous abordons notre dernière construction avec celle des réels. Il existe deux
constructions classiques de $\bR$~: celle par suites de Cauchy et celle
par coupures de Dedekind. Comme nous avons abordé la complétion dans le cadre
général dans le \cref{chp.topo}, nous allons étudier ici la construction par
les coupures de Dedekind, qui se base sur la structure ordonnée de $\bQ$.
Plutôt que de simplement donner la construction historique de
\cite{DedekindCut}, nous adoptons une construction adaptée de la complétion
idéale, présentée par exemple dans \cite{DBLP:books/daglib/0093287}.

La relation $\leq$, en effet, peut s'étendre de $\bN$ à $\bZ$ puis à $\bQ$,
conférant à $\bQ$ une structure de corps ordonné. Il nous faut donc commencer
par donner les définitions d'ordre sur $\bZ$ et $\bQ$~:

\begin{definition}
  On définit sur $\bZ$ la relation $\leq$ par~:
  \[a \leq b \defeq b - a \in \bN\]
  où $\bN$ est identifié à $\iota(\bN) \subseteq \bZ$.
\end{definition}

\begin{exercise}
  Montrer que $\leq$ est une relation d'ordre totale compatible avec l'addition,
  c'est-à-dire que pour tous $a\leq b$ et $c\leq d$, on a $a+c\leq b+d$.
  Montrer de plus que si $a \geq 0$ et $b\geq 0$ alors $a\times b \geq 0$.
  La structure $(\bZ,0,+,1,\times,\leq)$ est alors appelée un anneau
  totalement ordonné.
\end{exercise}

\begin{definition}
  On définit la partie de $\bQ$ suivante~:
  \[\bQ_+ \defeq \compr{\frac{p}{q}}{p,q\in \bN}\]
  On définit alors la relation $\leq$ sur $\bQ$ par~:
  \[r \leq s \defeq s - r \in \bQ_+\]
\end{definition}

\begin{exercise}
  Montrer que $(\bQ,0,+,1,\times,\leq)$ est aussi un anneau totalement
  ordonné. On dit alors que c'est un corps totalement ordonné.
\end{exercise}

De plus, toutes ces structures ordonnées ont la propriété d'être archimédiennes,
ce qui signifie que $\bZ$ et $\bQ$ vérifient la proposition
suivante~:
\[\forall k\in A, \exists n \in\bN, n \geq k\]

On veut désormais compléter $\bQ$, ce qui signifie moralement qu'on veut
supprimer tous les trous que contient cet ensemble. Dans le contexte des suites
de Cauchy, cela signifie qu'on veut qu'une suite de Cauchy, devant converger en
terme d'à quel point ses points se rapprochent, doit converger effectivement
vers un point donné. Dans le contexte de l'ordre, la notion naturelle de limite
est celle de borne supérieure ou inférieure. On peut voir cela par exemple avec
l'intervalle $\intOO{0}{1}$~: il ne contient ni $0$ ni $1$, mais on pense à
$1$ comme à un point limite de l'intervalle. Pour une suite croissante,
d'ailleurs, la limite coïncide avec la borne supérieure (dans $\bR$).

On souhaite donc ajouter à notre ensemble $\bQ$ toutes les bornes
supérieures de parties de $\bQ$. Une construction générale de théorie des
ordres existe pour cela~: c'est la complétion idéale. Elle consiste à étudier,
à la place d'un ensemble ordonné $(X,\leq)$, l'ensemble $\ideal(X)$ de ses
idéaux muni de l'inclusion. La construction des coupures de Dedekind est très
proche, mais considère au lieu des idéaux des paires $(L,U)$ d'un idéal et d'un
filtre. Dans notre cas, nous allons traiter le cas général de la complétion
idéale plutôt que la version historique des coupures de Dedekind.

Notre présentation sera un entre deux~: nous donnons une construction
similaire à celle de la complétion idéale pour les ordres partiels complets,
mais nous nous plaçons directement dans le cadre des ordres totaux. On parle
donc d'ordre total complet, et on considère une variante qui nous rapproche
directement de la propriété de la borne supérieure, qu'on veut vérifier dans
$\bR$.

\begin{definition}[Ordre total complet]
  On dit qu'un ensemble totalement ordonné $(X,\leq)$ est complet si toute
  partie non vide majorée admet une borne supérieure et toute partie minorée
  non vide admet une borne inférieure.
\end{definition}

Notre objectif est donc, à partir de $(\bQ,\leq)$, de construire
l'ensemble ordonné $(\bR,\leq)$ qui est complet et contient, en un sens,
$\bQ$. On veut donc \og ajouter les bornes supérieures\fg de $\bQ$,
les bornes inférieures en étant conséquences du fait que $x\mapsto -x$ envoie
les bornes supérieures sur les bornes inférieures et inversement (on le verra
après notre construction).

Partons d'une remarque essentielle~: les bornes supérieures existent toujours
quand on étudie des ensembles d'ensembles, grâce à l'union. Cela nous motive à
plonger $\bQ$ dans un ensemble d'ensembles. En l'occurrence, il existe une
façon assez naturelle d'associer une partie de $\bQ$ à un rationnel $r$~:
il suffit de considérer l'ensemble $\compr{r' \in \bQ}{r' \leq r}$. Cette
association tient compte de l'ordre, ce qui est souhaité vu notre objectif.

On veut donc caractériser les parties qui se comportent comme celle définie plus
haut. Une propriété évidente est qu'une telle partie est close par le bas, par
transitivité de l'ordre $\leq$. Une telle partie est aussi non vide, dans le cas
de $\bQ$, puisque $\bQ$ n'est pas minoré. Elle n'est pas non plus
l'ensemble tout entier. Ces conditions sont exactement celles définissant un
idéal dans un ensemble totalement ordonné (car la propriété à propos de $\opmt$
est alors automatiquement vérifiée), et ne pas être $\bQ$ signifie être un
idéal propre.

Cependant, on voit que deux idéaux peuvent sembler coder le même rationnel~:
$\{q < 2\}$ et $\{q \leq 2\}$ sont des idéaux et ont la même borne supérieure,
ce qui est un problème dans notre cas. D'un point de vue de la théorie des
ordres, on peut voir ces deux parties comme respectivement un ouvert et un
fermé, et il est donc nécessaire de choisir si l'on considère qu'un idéal
doit être ouvert ou fermé. On considère ici qu'un idéal est une partie ouverte,
ce qui signifie qu'un idéal ne doit pas avoir de maximum\footnote{si on avait
voulu définir un idéal comme une partie fermé, on aurait au contraire demandé
que si la borne supérieure existe alors elle appartient à l'idéal.}.

\begin{definition}[Complétion idéale]
  Soit $(X,\leq)$ un ensemble totalement ordonné. On appelle complétion idéale
  l'ensemble $\ideal(X)\subseteq \powerset(X)$ défini par le fait que
  $\varId\in \ideal(X)$ si et seulement si~:
  \begin{itemize}
  \item pour tout $x,y\in X$, si $x\leq y$ et $y \in \varId$ alors
    $x\in\varId$~;
  \item $\varId\neq \vide$~;
  \item $\varId\neq X$~;
  \item pour tout $x \in \varId$, il existe $y \in \varId$ tel que $x < y$.
  \end{itemize}
  muni de l'inclusion comme ordre.
\end{definition}

\noindent
On peut alors montrer que $(\ideal(X),\subseteq)$ est un ordre total complet.

\begin{proposition}
  Soit $(X,\leq)$ un ordre totalement ordonné non majoré. Alors
  $(\ideal(X),\subseteq)$ est un ordre total complet.
\end{proposition}

\begin{proof}
  On souhaite d'abord montrer que cet ordre est total. Soient
  $\varId,\varIdB\in\ideal(X)$, $\varId\neq \varIdB$. Par définition, on trouve
  $x\in \varId\diffsym \varIdB$. Par symétrie du rôle de
  $\varId$ et $\varIdB$, on suppose sans perte de généralité que $x\in \varId$.
  Alors pour tout $z\in \varIdB$, si $x \leq z$ on aurait $x\in \varIdB$ puisque
  $\varIdB$ est clos par le bas. Donc, comme $x\in \varId\setminus \varIdB$, on
  en déduit que $x \leq z$ est faux. Comme l'ordre est total, on en déduit que
  $z\leq x$, donc par clôture par le bas de $\varId$ que $z\in \varId$. En
  conclusion, $\varIdB\subseteq \varId$, donc l'ordre est total.

  Soit maintenant une partie $\varFam\subseteq \ideal(X)$ non vide majorée, on
  construit la borne supérieure de $\varFam$ en deux temps. On pose le majorant
  $\varId = \bigcup \varFam$. Vérifions que $\varId \in \ideal(X)$~:
  \begin{itemize}
  \item soient $x,y \in X$ tels que $x < y$ et $y \in \varId$. Par définition,
    on trouve $\varId_0 \in \varFam$ tel que $y \in \varId_0$, mais alors
    (puisque $\varId_0$ est un idéal) $x \in \varId_0$, donc $x \in \varId$.
  \item puisque la partie est non vide et que les idéaux sont supposés non
    vides, l'union $\varId$ est non vide.
  \item on a supposé que $\varFam$ est majorée~: on trouve donc $\varIdB$
    tel que pour tout $\varId_0 \in \varFam$, $\varId_0\subseteq \varIdB$,
    donc (comme l'union est la borne supérieure pour l'inclusion) on en déduit
    que $\varId\subseteq\varIdB$. Comme $\varIdB\neq X$, il existe un élément
    qui est hors de $\varIdB$ et qui doit alors être strictement supérieur à
    tous les éléments de $\varIdB$ (car $X$ est totalement ordonné). Cet élément
    est alors strictement supérieur à tous les éléments de $\varId$, donc
    $\varId\neq X$.
  \item soit $x \in \varId$. Par hypothèse, on trouve $\varId_0\in\varFam$ tel
    que $x \in \varId_0$, mais alors il existe $y \in \varId_0$ tel que
    $y > x$, donc il existe $y \in \varId$ tel que $y > x$.
  \end{itemize}
  Donc $\varId$ est bien un idéal, et c'est la borne supérieure de $\varFam$
  car c'est un majorant pour l'inclusion (en tant qu'union) et tout majorant
  de $\varFam$ pour l'inclusion contient en particulier $\bigcup\varFam$.

  \noindent
  Le cas des bornes inférieures est laissé en exercice.
\end{proof}

\begin{exercise}
  Montrer qu'une partie minorée non vide $\varFam$ possède une borne inférieure
  dans $(\ideal(X),\subseteq)$. \textit{Indication~: on pourra s'inspirer de
    la \cref{prop.treillis.complet.justInf}.}
\end{exercise}

\begin{exercise}\label{exo.ide.maj}
  Soit $(X,\leq)$ un ensemble totalement ordonné non majoré. Montrer que tout
  $\varId \in \ideal(X)$ est majoré dans $X$.
\end{exercise}

\begin{remark}
  On peut noter la différence de traitement de la borne supérieure de l'idéal
  entre l'intersection et l'union. La raison à cela peut s'interpréter de façon
  topologique~: nous imposons à nos idéaux d'être des ouverts, donc considérer
  l'union fonctionne sans heurt, tandis que l'intersection arbitraite demande
  de considérer l'intérieur.
\end{remark}

On veut aussi montrer que $\ideal(X)$ est un prolongement naturel, en un sens,
de $X$. Cela se traduit par deux faits~: on peut plonger $X$ dans $\ideal(X)$,
et $\ideal(X)$ est le plus petit ensemble contenant $X$ et étant stable par
sups de parties bornées non vides. Cependant, cela demande une dernière
condition technique, qui est la densité de l'ordre.

\begin{definition}
  Un ensemble ordonné $(X,\leq)$ est dit dense si pour tous $x,y \in X$ tels
  que $x < y$ il existe $z\in X$ tel que $x < z < y$.
\end{definition}

\begin{proposition}
  Soit $(X,\leq)$ un ensemble totalement ordonné dense non majoré ni minoré. La
  fonction
  \[\makeFun{\iota}{X}{\ideal(X)}{x}{\compr{y \in X}{y < x}}\]
  est une fonction strictement croissante.
\end{proposition}

\begin{proof}
  Montrons d'abord que la fonction est bien définie, c'est-à-dire que pour tout
  $x\in X$, $\iota(x)$ est bien un élément de $\ideal(X)$~:
  \begin{itemize}
  \item si $y,z\in X$ tels que $y\leq z$ et $z\in \iota(x)$, alors $z \leq x$
    donc $y\leq x$ par transitivité, donc $y\in \iota(x)$.
  \item $X$ n'est pas minoré, donc on trouve $y < x$, ce qui signifie que
    $y \in \iota(x)$, donc $\iota(x) \neq\vide$.
  \item $X$ n'est pas majoré, donc on trouve $y > x$, ce qui signifie que
    $y\notin\iota(x)$, donc $\iota(x)\neq X$.
  \item si $y \in \iota(x)$ alors $y < x$, et par densité de $X$ on en déduit
    qu'il existe $z\in X$ tel que $y < z < x$, c'est-à-dire $z \in \ideal(x)$
    tel que $z > y$.
  \end{itemize}
  On en déduit que $\iota$ est bien définie.

  Si $x < y$ alors pour tout $z \in \iota(x)$, $z < x$ donc par transitivité
  $z < y$, donc $z \in \iota(y)$. De plus, $y\notin\iota(x)$ et
  $y\in\iota(y)$ donc $\iota(x)\subsetneq\iota(y)$. $\iota$ est donc strictement
  croissante.
\end{proof}

\begin{remark}
  On peut utiliser ce résultat pour prouver que $\ideal(X)$ n'est pas majoré.
  En effet, s'il l'était, alors son majorant contiendrait tous les $\iota(x)$
  pour $x\in X$, donc tout $X$, ce qui n'est pas possible. En réalité, on
  pourrait simplement considérer l'ensemble des idéaux et non l'ensemble des
  idéaux propres (un idéal est propre s'il est différent de $X$), mais dans
  notre construction de $\bR$ cela introduirait un point à l'infini.
\end{remark}

\begin{lemma}\label{lem.id.cup}
  Soit $(X,\leq)$ un ensemble totalement ordonné dense non majoré ni minoré.
  Alors pour tout $\varId\in\ideal(X)$, on a l'égalité
  \[\varId = \bigcup_{x\in \varId} \iota(x)\]
\end{lemma}

\begin{proof}
  On prouve ce résultat par double inclusion~:
  \begin{itemize}
  \item pour tout $x\in \varId$, comme $\varId$ est clos par le bas,
    $\iota(x)\subseteq \varId$, donc
    $\displaystyle\bigcup_{x\in \varId}\iota(x)\subseteq \varId$.
  \item on sait que pour tout $x\in \iota(x)$ il existe $y \in \varId$ tel que
    $x < y$, donc $\displaystyle x\in \bigcup_{y\in \varId} \iota(y)$.\qedhere
  \end{itemize}
\end{proof}

\begin{theorem}[Propriété universelle de la complétion idéale]
  Soit $(X,\leq)$ un ensemble totalement ordonné dense non majoré ni mnioté.
  Soit $(Y,\leq)$ un ensemble totalement ordonné complet et
  $\makeFunName{f}{X}{Y}$ une fonction croissante préservant les bornes
  inférieures et supérieures. Il existe une unique fonction
  $\makeFunName{\relev{f}}{\ideal(X)}{Y}$ qui préserve les bornes supérieures et
  les bornes inférieures et telle que $f = \relev{f} \circ \iota$.
\end{theorem}

\begin{proof}
  On procède par analyse-synthèse~:
  \begin{itemize}
  \item si une telle fonction $\relev{f}$ existe, alors par définition, on sait
    que pour tout $x\in X$, $\relev{f}(\iota(x)) = f(x)$. De plus, par le
    \cref{lem.id.cup}, on sait que pour tout $\varId\in \iota(x)$,
    $\displaystyle \varId = \bigcup_{x\in \varId} \iota(x)$, donc par préservation
    des bornes supérieures on en déduit que
    \[\relev{f}(\varId) = \bigvee_{x\in \varId} f(x)\]
    ce qui définit de façon unique $\relev{f}$.
  \item Montrons que notre candidat $\relev{f}$ est une fonction croissante qui
    préserve les bornes supérieures et inférieures~:
    \begin{itemize}
    \item Tout d'abord, si
      $\varId\subseteq \varIdB$, comme la borne supérieure est croissante pour
      l'inclusion, on en déduit que $\opJn f(\varId) \leq \opJn f(\varIdB)$,
      d'où $\relev{f}(\varId)\leq \relev{f}(\varIdB)$.
      
    \item Soit $\varFam\subseteq \ideal(X)$ admettant une borne supérieure,
      qu'on note $\varId$. On a vu que cette borne supérieure est
      $\bigcup \varFam$, ce qui nous donne donc
      \[\relev{f}(\varId) = \opJn_{x \in \varId} f(x)\]
      on veut montrer que cette valeur est la borne supérieure de
      $\relev{f}(\varFam)$. Comme la borne supérieure est croissante vis à vis
      de l'inclusion, on sait déjà que $\relev{f}(\varId)$ est un majorant de
      $\relev{f}(\varFam)$.
      Soit $\varIdB$ un majorant de $\relev{f}(\varId)$, cela signifie que pour
      tout $\varId_0\in \varFam, \opJn_{x \in \varId_0} f(x) \leq \varIdB$, donc
      pour tout $x \in \varId, f(x) \leq \varIdB$.
    \item Le cas de la borne inférieure est laissé en exercice.
    \end{itemize}
  \end{itemize}
  Ainsi il existe une unique fonction $\relev{f}$ vérifiant les conditions
  imposées.
\end{proof}

\begin{exercise}
  Montrer que $\relev{f}$ préserve bien les bornes inférieures.
\end{exercise}

\begin{remark}
  On peut représenter la propriété précédente par le fait qu'il existe une
  unique flèche $\relev{f}$ faisant commuter le diagramme suivant~:
  \begin{center}
    \begin{tikzcd}
      X \ar[r,"f"]\ar[d,"\iota"] & Y\\
      \ideal(X) \ar[ur,dashed,"\relev{f}"]
    \end{tikzcd}
  \end{center}
\end{remark}

\begin{definition}[Ensemble des réels]
  On définit $(\bR,\leq)$ comme la complétion idéale de $(\bQ,\leq)$, qui est
  bien un ensemble totalement ordonné dense non majoré ni minoré.
\end{definition}

On définit finalement les deux lois $+$ et $\times$ sur $\bR$. Pour le cas de
l'addition, la construction est directe, mais la multiplication demande de
prendre en compte le signe des éléments.

\begin{definition}[Somme de réels]
  Soient $\varId,\varIdB\in\bR$, on définit la somme
  \[\varId+\varIdB\defeq \compr{x + y}{x \in \varId,y \in \varIdB}\]
  et l'opposé
  \[(-\varId) \defeq
  \compr{y}{\exists z,y < z \land (\forall x \in \varId, z < - x)}\]
\end{definition}

\begin{remark}
  Une autre façon de définir $(-\varId)$ est de considérer d'abord l'opposé
  $\compr{-x}{x \in \varId}$. Cet ensemble n'est pas un idéal car il est clos
  par le haut (prendre l'opposé inverse l'ordre)~: on souhaite donc considérer
  son complémentaire, mais le complémentaire est alors un fermé (car l'opposé
  d'un ouvert est ici un ouvert). On souhaite donc prendre l'intérieur de ce
  fermé, ce qui donne le même résultat.
\end{remark}

\begin{exercise}
  Vérifier que $(\bR,+,0)$ est un groupe abélien.
\end{exercise}

\begin{definition}[Produit de réels]
  Le produit de nombres réels $\varId,\varIdB > 0$ est défini par
  \[\varId\times_{\bR^*}\varIdB\defeq
  \compr{z \in \bQ}{\exists x \in \varId, \exists y \in \varIdB,
    (x > 0) \land (y > 0) \land (z < x \times y)}\]
  il est ensuite défini sur tous les réels par
  \[\varId \times \varIdB \defeq \begin{cases}
  0 \text{ si } \varId = 0 \lor \varIdB = 0 \\
  \varId \times_{\bR_+^*} \varIdB \text{ si } \varId > 0 \land \varIdB > 0\\
  - (\varId \times_{\bR_+^*} (-\varIdB)) \text{ si }
  \varId > 0 \land \varIdB < 0\\
  - ((-\varId) \times_{\bR_+^*} \varIdB) \text{ si }
  \varId < 0 \land \varIdB > 0\\
  - ((-\varId) \times_{\bR_+^*} (-\varIdB)) \text{ si }
  \varId < 0 \land \varIdB < 0
  \end{cases}\]
\end{definition}

\begin{exercise}
  Montrer que $(\bR,+,\times,\leq)$ est un corps totalement ordonné.
\end{exercise}

\begin{corollary}
  Toute partie non vide majorée (respectivement, minorée) de $\bR$ admet
  une borne supérieure (respectivement, inférieure).
\end{corollary}

\begin{remark}
  On aurait pu, au lieu de faire cette disjonction de cas, construire nos
  ensembles de nombres dans un ordre plus efficace en vu de la construction par
  complétion idéale. En construisant $\bN$, puis $\bQ_+$ (en
  inversant $\times$ dans $\bN\setminus\{0\}$), on peut directement considérer
  le demi-anneau $(\bR_+,+,\times,\leq)$ complet, puis le symétriser pour la
  loi $+$ pour obtenir un corps totalement ordonné complet.
\end{remark}

\begin{remark}
  L'autre construction de $\bR$ n'exige pas de construction
  \textit{ad hoc} comme la disjonction de cas précédente. Elle utilise la
  complétion par les suites de Cauchy~: il se trouve que si $(\varAA,+,\times)$
  est un anneau, alors les suites de Cauchy à valeurs dans $\varAA$ ont
  naturellement une structure d'anneau en considérant les opérations termes à
  termes. L'ensemble des suites de Cauchy de limite nulle est alors un idéal de
  cet anneau~: on peut donc directement construire l'anneau quotient. Cette
  construction de $\bR$ est plus directe et plus naturelle depuis $\bQ$, mais
  cette construction alternative est historiquement importante, et revêt aussi
  une grande importance dans le cas des mathématiques constructives, où les
  deux constructions de $\bR$ ne sont plus équivalentes (il est notamment
  possible dans ce cadre d'avoir une construction des réels de Dedekind qui
  donne un ensemble dénombrable \cite{BauerCountReal}).
\end{remark}
